% ==================================================
% Summation Operator / Discrete Antiderivative
% Discrete Calculus
% ==================================================

\subsection{Discrete Antiderivatives}
\label{sec:discrete-antiderivative}

The telescoping identity established in the previous section,
\[
\sum_{k=a}^{b-1} \Delta f(k) = f(b) - f(a),
\]
reveals that summation acts as an inverse operation to the difference
operator, much like integration acts as an inverse to differentiation
in ordinary calculus. This motivates the following definition.

\begin{definition}[Discrete Antiderivative]
\label{def:discrete-antiderivative}
Let \(f : \mathbb{Z} \to \mathbb{R}\) be a discrete function.
A function \(F : \mathbb{Z} \to \mathbb{R}\) is called a
\emph{discrete antiderivative} (or \emph{discrete indefinite sum}) of \(f\) if
\[
\Delta F(n) = f(n),
\]
that is, \(F(n+1) - F(n) = f(n)\) for all \(n \in \mathbb{Z}\).
\end{definition}

\begin{remark*}[Standard quantified statement]
\(\forall f,F : \mathbb{Z}\to\mathbb{R}:\;
F \text{ is a discrete antiderivative of } f
\iff
\forall n \in \mathbb{Z},\; F(n+1) - F(n) = f(n)\).
\end{remark*}

\begin{remark*}[Interpretation]
A discrete antiderivative is a function whose one-step increments reproduce
the given summand.
\end{remark*}

\begin{remark*}[Analogy]
A discrete antiderivative plays the same role in discrete calculus that an
antiderivative plays in ordinary calculus: it reverses the fundamental
differentiation-like operator of the theory. Instead of reversing
differentiation, it reverses the forward difference \(\Delta\).
\end{remark*}

\begin{dependencies}
\begin{itemize}
  \item \hyperref[def:forward-difference]{Forward Difference Operator}.
\end{itemize}
\end{dependencies}

\subsection*{Examples}

\textbf{Example 1.} Let \(f(n) = 1\). Take \(F(n) = n\). Then
\[
\Delta F(n) = (n+1) - n = 1 = f(n). \checkmark
\]

\textbf{Example 2.} Let \(f(n) = n\). Take \(F(n) = \tfrac{n(n-1)}{2}\). Then
\[
\Delta F(n)
= \frac{(n+1)n}{2} - \frac{n(n-1)}{2}
= \frac{n(n+1-n+1)}{2}
= n = f(n). \checkmark
\]

\begin{remark*}[Connection to falling factorials]
Both examples above are instances of a general pattern: the discrete
antiderivative of the falling factorial \(n^{\underline{k}}\) is
\(\tfrac{1}{k+1}n^{\underline{k+1}}\), exactly mirroring the antiderivative
rule for ordinary powers. This is developed in the Polynomial Basis section.
\end{remark*}

\subsection*{Constant of Integration}

Discrete antiderivatives are not unique.

\begin{proposition}[Constant of integration]
\label{prop:constant-of-integration}
If \(F\) is a discrete antiderivative of \(f\), then so is \(F + C\)
for any constant \(C \in \mathbb{R}\).
\end{proposition}

\begin{remark*}[Fully quantified form]
\(\forall C \in \mathbb{R}:\;
\Delta F = f \Rightarrow \Delta(F+C) = f\).
\end{remark*}

\begin{remark*}[Interpretation]
Adding a constant changes the level of a discrete antiderivative but not its
one-step increments.
\end{remark*}

\begin{dependencies}
\begin{itemize}
  \item \hyperref[def:discrete-antiderivative]{Discrete Antiderivative}.
  \item \hyperref[def:forward-difference]{Forward Difference Operator}.
\end{itemize}
\end{dependencies}

\begin{proof}
\[
\Delta(F+C)(n) = (F+C)(n+1) - (F+C)(n)
= F(n+1)-F(n) = \Delta F(n) = f(n). \qedhere
\]
\end{proof}
